\chapter[Il Nobile Ottuplice Sentiero]{L'insegnamento sul Nobile Ottuplice Sentiero}

\firstline{Ayam-eva ariyo aṭṭhaṅgiko maggo}

\begin{leader}
  [Handa mayaṁ ariyaṭṭhaṅgika-magga-pāṭham bhaṇāmase]
\end{leader}

\begin{english}
  [Cantiamo ora il sentiero del Nobile Ottuplice Sentiero]
\end{english}

Ayam-eva a꜕riyo aṭṭha꜓ṅgi꜕ko maggo

\begin{english}
  Questo è il Nobile Ottuplice Sentiero,
\end{english}

Se꜓yyathī꜓daṁ

\begin{english}
  Che è come segue:
\end{english}

Sa꜓mmā-diṭṭhi

\begin{english}
  Retta Visione,
\end{english}

Sa꜓mmā-sa꜓ṅka꜕ppo

\begin{english}
  Retta Intenzione,
\end{english}

Sa꜓mmā-vācā

\begin{english}
  Retta Parola,
\end{english}

Sa꜓mmā-kammanto

\begin{english}
  Retta Azione,
\end{english}

Sa꜓mmā-ājīvo

\begin{english}
  Retto Sostentamento,
\end{english}

Sa꜓mmā-vā꜕yāmo

\begin{english}
  Retto Sforzo,
\end{english}

\ifaivedition
\clearpage
\fi

Sa꜓mmā-sa꜕ti

\begin{english}
  Retta Consapevolezza,
\end{english}

Sa꜓mmā-sa꜕mādhi

\begin{english}
  Retta Concentrazione.
\end{english}

Ka꜕tamā ca bhi꜓kkh꜕ave sammā-diṭṭhi

\begin{english}
  E cos'è, bhikkhu, la Retta Visione?
\end{english}

Yaṁ kho bhi꜓kkh꜕ave dukkhe ñāṇaṁ

\begin{english}
  La conoscenza della sofferenza;
\end{english}

Dukkha-sa꜕mu꜕daye ñāṇaṁ

\begin{english}
  La conoscenza dell'origine della sofferenza;
\end{english}

Dukkha-ni꜓rodhe ñāṇaṁ

\begin{english}
  La conoscenza della cessazione della sofferenza;
\end{english}

Dukkha-ni꜓rodha-gāmi꜓ni꜓yā pa꜕ṭipa꜕dāya ñāṇaṁ

\begin{english}
  La conoscenza del sentiero che porta alla cessazione della sofferenza:
\end{english}

A꜕yaṁ vuccati bhi꜓kkh꜕ave sa꜓mmā-diṭṭhi

\begin{english}
  Questa, bhikkhu, è chiamata Retta Visione.
\end{english}

Katamo ca bhi꜓kkh꜕ave sammā-sa꜓ṅka꜕ppo

\begin{english}
  E cos'è, bhikkhu, la Retta Intenzione?
\end{english}

\ifaivedition
\clearpage
\fi

Nekkhamma-sa꜓ṅka꜕ppo

\begin{english}
  L'intenzione di rinuncia;
\end{english}

A꜕byāpāda-sa꜓ṅka꜕ppo

\begin{english}
  L'intenzione di non-malevolenza;
\end{english}

A꜕vihiṁsā-sa꜓ṅka꜕ppo

\begin{english}
  L'intenzione di non-crudeltà:
\end{english}

Ayaṁ vuccati bhi꜓kkh꜕ave sa꜓mmā-sa꜓ṅka꜕ppo

\begin{english}
  Questa, bhikkhu, è chiamata Retta Intenzione.
\end{english}

Katamā ca bhi꜓kkh꜕ave sa꜓mmā-vācā

\begin{english}
  E cos'è, bhikkhu, la Retta Parola?
\end{english}

Musā-vādā vera꜓ma꜕ṇī

\begin{english}
  Astenersi dalla parola falsa;
\end{english}

Pisuṇāya vācāya vera꜓ma꜕ṇī

\begin{english}
  Astenersi dalla parola maliziosa;
\end{english}

Pharusāya vācāya vera꜓ma꜕ṇī

\begin{english}
  Astenersi dalla parola aspra;
\end{english}

Sa꜓mphappa꜕lāpā vera꜓ma꜕ṇī.

\begin{english}
  Astenersi dal chiacchiericcio futile:
\end{english}

\ifaivedition
\clearpage
\fi

Ayaṁ vuccati bhi꜓kkh꜕ave sa꜓mmā-vācā

\begin{english}
  Questa, bhikkhu, è chiamata Retta Parola.
\end{english}

Katamo ca bhi꜓kkh꜕ave sa꜓mmā-kammanto

\begin{english}
  E cos'è, bhikkhu, la Retta Azione?
\end{english}

Pāṇāti꜕pātā vera꜓ma꜕ṇī

\begin{english}
  Astenersi dall'uccidere esseri viventi;
\end{english}

A꜕dinnādānā vera꜓ma꜕ṇī

\begin{english}
  Astenersi dal prendere ciò che non è dato;
\end{english}

Kāmesu꜕ micchā꜓cārā vera꜓ma꜕ṇī

\begin{english}
  Astenersi dalla cattiva condotta sessuale:
\end{english}

Ayaṁ vuccati bhi꜓kkh꜕ave sa꜓mmā-kammanto

\begin{english}
  Questa, bhikkhu, è chiamata Retta Azione.
\end{english}

Katamo ca bhi꜓kkha꜕ve sa꜓mmā-ājīvo

\begin{english}
  E cos'è, bhikkhu, il Retto Sostentamento?
\end{english}

Idha bhi꜓kkh꜕ave a꜕riya-sā꜓va꜕ko micchā-ājīvaṁ pa꜕hāya sammā-ājī꜓vena jīvi꜓taṁ ka꜕ppeti

\begin{english}
  Qui, bhikkhu, un Nobile Discepolo, avendo abbandonato un errato sostentamento, si guadagna da vivere con un retto sostentamento:
\end{english}

Ayaṁ vuccati bhi꜓kkh꜕ave sa꜓mmā-ājīvo

\begin{english}
  Questo, bhikkhu, è chiamato Retto Sostentamento.
\end{english}

Katamo ca bhi꜓kkh꜕ave sa꜓mmā-vāyāmo

\begin{english}
  E cos'è, bhikkhu, il Retto Sforzo?
\end{english}

Idha bhi꜓kkh꜕ave bhikkhu a꜕nuppannānaṁ pāpa꜕kānaṁ a꜕ku꜕salānaṁ dhammānaṁ anuppādāya

\begin{english}
  Qui, bhikkhu, un bhikkhu risveglia lo zelo per il non-sorgere di stati malvagi e non salutari non ancora sorti;
\end{english}

Chandaṁ ja꜕neti vāyama꜓ti vī꜓ri꜓yaṁ ārabha꜕ti ci꜕ttaṁ pa꜕ggaṇhā꜓ti pa꜕daha꜕ti

\begin{english}
  Egli si sforza, suscita energia, esercita la sua mente e lotta.
\end{english}

U꜕ppannānaṁ pāpa꜕kānaṁ a꜕ku꜕salānaṁ dhammānaṁ pa꜕hānāya

\begin{english}
  Egli risveglia lo zelo per l'abbandono degli stati malvagi e non salutari sorti;
\end{english}

Chandaṁ ja꜕neti vāyama꜓ti vī꜓ri꜓yaṁ ārabha꜕ti ci꜕ttaṁ pa꜕ggaṇhā꜓ti pa꜕daha꜕ti

\begin{english}
  Egli si sforza, suscita energia, esercita la sua mente e lotta.
\end{english}

Anuppannānaṁ ku꜕salānaṁ dhammānaṁ u꜕ppādāya

\begin{english}
  Egli risveglia lo zelo per il sorgere di stati salutari non ancora sorti;
\end{english}

Chandaṁ ja꜕neti vāyama꜓ti vī꜓ri꜓yaṁ ārabha꜕ti ci꜕ttaṁ pa꜕ggaṇhā꜓ti pa꜕daha꜕ti

\begin{english}
  Egli si sforza, suscita energia, esercita la sua mente e lotta.
\end{english}

\ifaivedition
\clearpage
\fi

U꜕ppannānaṁ ku꜕salānaṁ dhammānaṁ ṭh꜓iti꜕yā a꜕sa꜕mmosāya bh꜓iyyobhāvāya vepu꜕llāya bhāva꜓nāya pāri꜓pū꜕riyā

\begin{english}
  Egli risveglia lo zelo per la continuazione, la non-scomparsa, il rafforzamento, l'aumento e il compimento attraverso lo sviluppo degli stati salutari sorti;
\end{english}

Chandaṁ ja꜕neti vāyama꜓ti vī꜓ri꜓yaṁ ārabha꜕ti ci꜕ttaṁ pa꜕ggaṇhā꜓ti pa꜕daha꜕ti

\begin{english}
  Egli si sforza, suscita energia, esercita la sua mente e lotta:
\end{english}

Ayaṁ vuccati bhi꜓kkh꜕ave sa꜓mmā-vāyāmo

\begin{english}
  Questo, bhikkhu, è chiamato Retto Sforzo.
\end{english}

Katamā ca bhi꜓kkh꜕ave sa꜓mmā-sa꜕ti

\begin{english}
  E cos'è, bhikkhu, la Retta Consapevolezza?
\end{english}

Idha bhi꜓kkh꜕ave bhikkhu kāye kāyānupa꜕ssī vi꜓ha꜕rati

\begin{english}
  Qui, bhikkhu, un bhikkhu dimora contemplando il corpo come corpo,
\end{english}

Ātāpī sa꜓mpa꜕jāno sa꜕timā

\begin{english}
  Ardente, pienamente consapevole e attento,
\end{english}

Vi꜓neyya loke a꜕bhijjhā-domanassaṁ

\begin{english}
  Avendo messo da parte bramosia e afflizione per il mondo;
\end{english}

Veda꜕nāsu꜕ veda꜕nānu꜓pa꜕ssī vi꜓ha꜕rati

\begin{english}
  Egli dimora contemplando le sensazioni come sensazioni,
\end{english}

\ifaivedition
\clearpage
\fi

Ātāpī sa꜓mpa꜕jāno sa꜕timā

\begin{english}
  Ardente, pienamente consapevole e attento,
\end{english}

Vi꜓neyya loke a꜕bhijjhā-domanassaṁ

\begin{english}
  Avendo messo da parte bramosia e afflizione per il mondo;
\end{english}

Ci꜕tte ci꜕ttānu꜓pa꜕ssī vi꜓ha꜕rati

\begin{english}
  Egli dimora contemplando la mente come mente,
\end{english}

Ātāpī sa꜓mpa꜕jāno sa꜕timā

\begin{english}
  Ardente, pienamente consapevole e attento,
\end{english}

Vi꜓neyya loke a꜕bhijjhā-domanassaṁ

\begin{english}
  Avendo messo da parte bramosia e afflizione per il mondo;
\end{english}

Dhammesu꜕ dhammānu꜓pa꜕ssī vi꜓ha꜕rati

\begin{english}
  Egli dimora contemplando gli oggetti mentali come oggetti mentali,
\end{english}

Ātāpī sa꜓mpa꜕jāno sa꜕timā

\begin{english}
  Ardente, pienamente consapevole e attento,
\end{english}

Vi꜓neyya loke a꜕bhijjhā-domanassaṁ

\begin{english}
  Avendo messo da parte bramosia e afflizione per il mondo:
\end{english}

Ayaṁ vuccati bhi꜓kkh꜕ave sa꜓mmā-sa꜕ti

\begin{english}
  Questa, bhikkhu, è chiamata Retta Consapevolezza.
\end{english}

\ifaivedition
\clearpage
\fi

Katamo ca bhi꜓kkh꜕ave sa꜓mmā-sa꜕mādhi

\begin{english}
  E cos'è, bhikkhu, la Retta Concentrazione?
\end{english}

Idha bhi꜓kkh꜕ave bhikkhu

\begin{english}
  Qui, bhikkhu, un bhikkhu,
\end{english}

Vivicc'eva kāmehi

\begin{english}
  Piuttosto appartato dai piaceri sensuali,
\end{english}

Vivicca a꜕ku꜕sa꜕lehi dh꜕ammehi

\begin{english}
  Appartato dagli stati non salutari,
\end{english}

Sa꜕vi꜓ta꜕kkaṁ sa꜕vi꜓cāraṁ viveka꜕-jaṁ pīti꜕-sukhaṁ pa꜕ṭhamaṁ jhānaṁ upasa꜓mpajja vi꜓ha꜕rati

\begin{english}
  Entra e dimora nel primo jhāna — accompagnato da pensiero applicato e sostenuto, con rapimento e piacere nati dal ritiro.
\end{english}

Vi꜓takka-vicārānaṁ vūpa꜕samā

\begin{english}
  Con la quiete del pensiero applicato e sostenuto,
\end{english}

Ajjhattaṁ sa꜓mpa꜕sādanaṁ ceta꜕so ekodi꜓bhāvaṁ avi꜓ta꜕kkaṁ aviInfocāraṁ sa꜕mādhi꜓-jaṁ pīti꜕-sukhaṁ du꜕tiyaṁ jhānaṁ upasa꜓mpa꜕jja vi꜓ha꜕rati

\begin{english}
  Egli entra e dimora nel secondo jhāna — accompagnato da fiducia in sé e unicità della mente, senza pensiero applicato e sostenuto, con rapimento e piacere nati dalla concentrazione.
\end{english}

\clearpage

Pītiyā ca꜕ vi꜓rāgā

\begin{english}
  Con lo svanire anche del rapimento
\end{english}

U꜕pekkhako ca vi꜓ha꜕rati

\begin{english}
  Egli dimora in equanimità,
\end{english}

Sa꜕to ca꜕ sa꜓mpa꜕jāno

\begin{english}
  Consapevole e pienamente cosciente,
\end{english}

Su꜕khañca kāyena pa꜕ṭisa꜓ṁvedeti

\begin{english}
  Provando ancora piacere con il corpo,
\end{english}

Yaṁ taṁ a꜕riyā āci꜕kkhanti u꜕pekkha꜓ko sa꜕timā su꜕kha-vi꜓hā꜕rī'ti tatiyaṁ jhānaṁ u꜕pasa꜓mpa꜕jja vi꜓ha꜕rati

\begin{english}
  Egli entra e dimora nel terzo jhāna — riguardo al quale i Nobili dichiarano: 'Egli ha una dimora piacevole, con equanimità ed è consapevole.'
\end{english}

Sukhassa ca꜕ pahānā

\begin{english}
  Con l'abbandono del piacere
\end{english}

Dukkhassa ca꜕ pahānā

\begin{english}
  E l'abbandono del dolore,
\end{english}

Pu꜕bb'eva somanassa꜕-domanassā꜓naṁ a꜕tthaṅga꜕mā

\begin{english}
  Con la precedente scomparsa di gioia e afflizione,
\end{english}

Adukkham-asu꜕khaṁ u꜕pekkhā-sa꜕ti-pā꜕ri꜓su꜕ddhiṁ ca꜕tutthaṁ jhānaṁ u꜕pasa꜓mpa꜕jja vi꜓ha꜕rati

\begin{english}
  Egli entra e dimora nel quarto jhāna — accompagnato né da dolore né da piacere, e purezza della consapevolezza dovuta all'equanimità:
\end{english}

Ayaṁ vuccati bhi꜓kkh꜕ave sa꜓mmā-sa꜕mādhi

\begin{english}
  Questo, bhikkhu, è chiamato Retta Concentrazione.
\end{english}

Ayam-eva a꜕riyo aṭṭha꜓ṅgi꜕ko maggo

\begin{english}
  Questo è il Nobile Ottuplice Sentiero.
\end{english}